\subsection{从语言模型到 Agent：阶段不同，瓶颈也不同}

开场五页把整讲压缩成一张研发地图。前两页提出“这些模型到底怎样训练”的总问题；第三页把 pretraining、reasoning RL 与经典 post-training/RLHF 并列；第四页指出 architecture 和 loss 并不是实践中的全部；第五页再把 prompting 与 finetuning 放到“已有模型如何专门化”的支线上。这里最重要的不是记住某个美元数字，而是看出资源形态的迁移：预训练吃海量 token 与长期计算，后训练吃高质量任务与可靠评估，reasoning RL 又额外依赖可交互环境和防作弊 verifier。

\lecturefigure{slides-images/slide-001.png}{官方 slide 1：课程题目、讲者与公开资料边界}
\lecturefigure{slides-images/slide-002.png}{官方 slide 2：从 LLM 普及转向“模型如何训练”}
\lecturefigure{slides-images/slide-003.png}{官方 slide 3：预训练、推理强化学习与经典后训练的数量级对照}
\lecturefigure{slides-images/slide-004.png}{官方 slide 4：训练流水线从架构/损失扩展到数据、评估与系统}
\lecturefigure{slides-images/slide-005.png}{官方 slide 5：prompting 与 finetuning 的专门化路径}

\begin{knowledgebox}{读图：先比较“瓶颈列”，不要先比较价格}
slide 3 的 Data、Time、Compute cost 和 Bottleneck 是同一套比较维度。预训练的数据量以十万亿 token 计，训练时间以月计；经典后训练的数据量小几个数量级，却更依赖样本质量和多轴评估；reasoning RL 的难点不是只把训练跑起来，而是构造不会被策略钻空子的环境。slide 5 进一步说明 prompting 几乎没有训练成本，但其瓶颈仍然是评估，因为没有任务集就无法判断提示词是否真的改善了行为。课程因此把“数据、评估、系统”放在贯穿所有阶段的位置。
\end{knowledgebox}

\begin{warningbox}{数量级不是报价单}
\textbf{课堂提示：}讲者在 00:00--00:09 多次说明这些数字来自 2025 年公开的 Kimi、LLaMA、DeepSeek 等项目，只用于建立 order-of-magnitude intuition。不同模型规模、硬件租赁方式、利用率和数据采购策略会显著改变实际成本，不能把“$>10M$”或“$>1M$”直接当成项目预算。真正稳定的结论是：阶段改变后，主要稀缺资源也随之改变。
\end{warningbox}

这组图还纠正了一个常见误区：Agent 能力不是在基础模型上“再加一个框架”就自动出现。预训练提供知识和一般模式，后训练塑造交互行为，reasoning RL 强化可验证推理，工具和环境训练再把文本决策变成有副作用的动作。任何一段缺少对应评估，下一段都可能把错误放大。因此全课后续的内容虽然跨越数据、RL、GPU 与并行系统，实际上都在回答同一个端到端问题。

